\code{ThreadTracking} is used to configure the handling of thread events for its
associated Process.

\begin{apient}
static void setDefaultTrackThreads(bool b)
\end{apient}

\apidesc{
Sets the default handling mechanism for thread events across all Process objects
to interrupt-driven (b = true) or polling-driven (b = false).
}

\begin{apient}
static bool getDefaultTrackThreads()
\end{apient}

\apidesc{
Returns the current default handling mechanism for thread events across all
Process objects. A return value of true indicates interrupt-driven while
false indicates polling-driven.
}

\begin{apient}
bool setTrackThreads(bool b) const
\end{apient}

\apidesc{
Sets the thread event handling mechanism for the associated Process object to
interrupt-driven (b = true) or polling-driven (b = false).
Returns true on success and false on failure.
}

\begin{apient}
bool getTrackThreads() const
\end{apient}

\apidesc{
Returns the current thread event handling mechanism for the associated Process
object. A return value of true indicates interrupt-driven while false
indicates polling-driven.
}

\begin{apient}
bool refreshThreads()
\end{apient}

\apidesc{
Manually polls for queued thread events to handle.
Returns true on success and false on failure.
}

